Tables are among the most pervasive formats for structured information, underpinning enterprise analytics, scientific reporting, and everyday data exchange. Yet reasoning over them with large language models (LLMs) is fundamentally difficult: models consume one-dimensional token sequences, whereas a table's meaning is carried jointly by its rows, columns, headers, and the relationships among cells. The difficulty is compounded by scale. Current methods are demonstrated on small, clean tables (often only a handful of columns and a few dozen rows), whereas the data that matters in practice lives in enterprise databases whose schemas are large, cryptic, and obsolescent and whose tables are far too big to fit in a model's context. Closing this gap between demonstration and deployment, making tabular understanding work at \emph{industry scale}, is the motivation of this work. 

This report surveys how the field bridges the representation mismatch and organizes it around two complementary pillars that together address that gap: \emph{schema understanding}, the mapping of natural-language intent onto a table's structural blueprint, and \emph{data understanding}, the problem of reasoning faithfully over table content, whose canonical task is Table Question Answering. Because how a table is represented to a model (serialization, data schema, image, a specialized table encoder, or a graph) bears most directly on reasoning over content, we treat representation within the data-understanding pillar, singling out graph representations as especially promising for making table structure explicit and addressable yet so far demonstrated only on small or document-centric data, whose extension to large tables we take as a motivating direction. We synthesize the literature under the two pillars, tracing a consistent shift from single-pass methods toward multi-step pipelines that decompose, retrieve, verify, and correct, and arguing that the binding constraints are increasingly architectural rather than a matter of raw model capacity.

Against this backdrop we present our contribution to the schema-understanding pillar, Reflect-SQL~\cite{jain2027reflectsql}, a self-reflective Text-to-SQL system that constructs a semantic knowledge base over the schema, retrieves relevant elements iteratively, and synthesizes queries under syntactic and multi-dimensional semantic validation, scoring correctness with an LLM-as-a-judge and without a gold reference. By grounding reflection in the checkable signals the SQL setting affords, Reflect-SQL reaches $72.03\%$ execution accuracy on the industry-scale BIRD benchmark.

Finally, we lay out an extensive research agenda across both pillars: for schema understanding, continually evolving knowledge bases, retrieval that scales by addressing pre-computed structure rather than searching at query time, cross-database join reasoning, and reliable yet inexpensive reflection; and for data understanding, our next direction, representing tabular data at scale as a graph so that evidence retrieval becomes addressing rather than search, with robustness on complex table structures and higher-order reasoning as the targets and a large-scale Table Question Answering benchmark as the instrument to measure them. The unifying thesis is that progress in tabular data understanding depends less on ever-larger models than on the architecture that supplies them with the right structure and the right evidence.
